\section{The Eighth Victory of Hercules Over the Mares of Diomedes}

Hercules is sent against another tyrant, Diomedes, who owns man-eating horses in barbarous Thrace, as king of the warlike Bistonians, descended from Mars. Since the myth teaches us that these horses to be fetched are mares, we can identify them as the female psychological side of the male, that is to say, the region of the subconscious psyche, where the Anima dwells, unawakened, and therefore unable to unify with the male psyche and heal its division. Diomedes is the apex of this condition, and represents the false Ego, which is oh-so-certain that it can control the subconscious forces of the mind, riding them to tyranny, and feeding their inordinate passion on the flesh of the innocent.

This condition represents a wide awake nightmare\footnote{\url{https://en.wikipedia.org/wiki/The_Nightmare}}, and the visual embodiment of such can be pictured here. Moderns are fond today of reinstating The Matriarchy. This represents a re-descent into the group, collective trajectory of mankind. It is of no more use to return to Matriarchy, and perhaps a great deal less, than it would be to reinstitute earlier phases of what is called ``the Patriarchy'': the Patriarchy was actually a step up in the struggle of man, and created the preconditions for the Age of Courtly Love, for it is only in eternal and infinite Love that man and woman can find qualitative equality. It was Love that summoned Abraham out of the tribes of Chaldea, that built the limes and roads of ancient Roma which lead to the sacred seven hills, and that raised the towers of the Gothic arches over Europe. Christ Himself said He came to earth not to bring peace, but a sword. To quote an old Internet acquaintance, when the man puts on his hat, and comes around, there is conflict. It is a difficult truth to face, but wars, plagues, devastation, and human suffering (the four horseman of the Apocalypse), are servants to the Living Word of God, who is brighter than the Sun and has a sword coming out of His mouth. This is not the Jesus that the Episcopal Church in America (our unofficial national liberal denomination along with Zionism) likes to paint from the pulpit every Sunday.

Because mankind is ``behind the level of the Times'' (unequal to the Cosmic Law's progression), there are periodic upheavals and convulsions on the earth. This is why Jesus comforted His followers: ``Fear not, for greater am I that are in you, than He that is in the world''. The prince of this world is Satan, and he is a deputy of Cosmic Law, fulfilling a very specific role in the economy of the fallen and fractured universe, ruling man with fear, hunger, and libido dominandi. It is he that will drag man, willy-nilly, collectively toward a final confrontation with the truth about human nature. Do not mistake this Zeitgeist for the Holy Spirit.

Unintegrated man cannot reassimilate to God and overcome the fractured condition of the Fall because he has lost his sleeping queen. He cannot respond in freedom towards God -- he is alienated, not just from God, but from the truth, from his true self, and from Reality. He is a thrall to Satan. Every single one of his one-sided reactions flow out of this fundamental aversion to coming to grips with Love. This is expressed poetically in Francis Thompson's \emph{The Hound of Heaven}\footnote{\url{http://www.maryhillmuseum.org/2013/wp-content/uploads/2013/03/hound\_of\_heaven\_flyer.pdf}}:

\begin{quotex}
\begin{verse}
Come then, ye other children, Nature's---share\\
With me' (said I) `your delicate fellowship;\\
Let me greet you lip to lip,\\
Let me twine with you caresses,\\
Wantoning With our Lady-Mother's vagrant tresses,\\
Banqueting With her in her wind-walled palace,\\
Underneath her azured daïs,\\
Quaffing, as your taintless way is,\\
From a chalice\\
Lucent-weeping out of the dayspring.'\\
So it was done:\\
I in their delicate fellowship was one---\\
Drew the bolt of Nature's secrecies.\\
I knew all the swift importings\\
On the wilful face of skies;\\
I knew how the clouds arise\\
Spumèd of the wild sea-snortings;\\
All that's born or dies\\
Rose and drooped with; made them shapers\\
Of mine own moods, or wailful or divine;\\
With them joyed and was bereaven.\\
I was heavy with the even,\\
When she lit her glimmering tapers\\
Round the day's dead sanctities.\\
I laughed in the morning's eyes.\\
I triumphed and I saddened with all weather,\\
Heaven and I wept together,\\
And its sweet tears were salt with mortal mine;\\
Against the red throb of its sunset-heart\\
I laid my own to beat,\\
And share commingling heat;\\
But not by that, by that, was eased my human smart.\\
In vain my tears were wet on Heaven's gray cheek.\\
For ah! we know not what each other says,\\
These things and I; in sound I speak---\\
Their sound is but their stir, they speak by silences.\\
Nature, poor stepdame, cannot slake my drouth;\\
Let her, if she would owe me,\\
Drop yon blue bosom-veil of sky, and show me\\
The breasts o' her tenderness:\\
Never did any milk of hers once bless\\
My thirsting mouth.
\end{verse}
\end{quotex}


If you were told that every single one of ``the world's problems'' has its ugly root in the darkest recesses of your own soul, would you give it serious energy and investigation? Or would you flee into deeper projections of your shadow? Some have even decided to join the shadow, and meld themselves as much as possible into that projection, to deepen the illusion past repair in this life, saving the extraordinary grace of God.''

\begin{quotex}
``You are as spiritual as you desire to be, that is all.'' They were somewhat annoyed at the abruptness of his words, and turned away. At once he spoke to them in a loving tone. ``My dear children, I said your spirituality was what you wish it to be so that you would understand that your spirituality is entirely in proportion to your good will. Then enter into yourselves: don't ask other about your progress. Examine your good will, and from that alone you will discover the measure of your spirituality.''
\flright{\textsc{John Ruysbroeck}, speaking to pilgrims, in \emph{European Mystics}, by Rudolf Steiner, p44.}
\end{quotex}


It doesn't matter what happens to the world, if you lose your own soul. This Anima is asleep deep within, the pearl of great price, an analogy for the Kingdom of Heaven, which is within you. It is Hercules, the deputy of the true self from the human aspect and the true self from the aspect of his immortal descent, who must tame the Mares of Diomedes. He and only he can can bring them to heel. This is part of the process of purification proceeding illumination and theosis. The personal unconscious has to be overcome before the subconscious can be reintegrated, and Hercules embodies a more or less purified consciousness who can face the subconscious.

Do not be deceived. The Mares of Diomedes do not feed on other horses. They feed on human flesh, and Hercules will turn them against the Satanic ego that is so hubristic as to think that they can wield them as a weapon against others. This motif, for those who are too modern for Greek myths, can be found in the Game of Thrones\footnote{\url{https://www.youtube.com/watch?v=2tuLYG9sv68}} as well. One way or the other, as we will see below, the un-purified personal consciousness will be unable to tame, or even keep locked up, the forces of the subconscious. Carl Jung remarked that men over 40 are either religious, or they are neurotic. Confucius gave the age of 40 as important, holding out hope for anyone that was straying, until they reached at least that age. A man ought to have a few things figured out by then. There is an old African proverb as well : ``After the age of forty, a man is responsible for his own face''. But for those who do make the choice of some higher path, and attain personal results, and embark on the labors, there is great reward and consolation for their struggle and sacrifice. They are able to begin the real work. We shouldn't think it is just a matter of ``getting saved'', or ``getting our act together philosophically''. There is a point beyond the point.

The Mares submit to Hercules (naturally), but he (for the time being) leaves them in the hands of his companion Abderis. One cannot try to tame the subconscious, or face it, or even behold it, with anything less than the highest sum of all that is true, beautiful, and best in one's self (the Soul), and so this leads to the death of his companion. The lesson here is that nothing can be held back, that full commitment is necessary, in order to journey into the shadow realm. Hercules turns back to face Diomedes, and his companion dies. One could also say that the old personality has to ``die'' in order for the Soul to complete the Labor.

Hercules is always getting sidetracked, albeit in an understandable and sympathetic way. There is an inevitable human sanguinity about him, whether using a river to wash out the stables, or going berserk and killing everyone he can get his hands on. We have seen this several times in the myths. People get hurt when this happens. Inevitably. Perhaps the best lesson one could draw from this is to not get sidetracked. Have you considered that wavering on the goal, although perhaps not fatal to yourself in the long run, may very well prove fatal to important parts of yourself, and incur karmic debt? Or to other, important people in your life story? Hercules will have to found a town in honor of his friend, which is another sidetrack, of a more forgivable kind. Although he is always ``in the right'' and ``doing something good'', it is not always, in every exactitude, the \emph{unum necessarium}. Be careful about pursuing sideline goals, and postponing your spiritual journey.

It is as if the entire story of Creation, all of written and unwritten human history, stretching back the flickering fire light of caves or even the time of proto-genesis, is a progressive tightening of focus around the ever-more-definable question of \emph{What does it mean to be Human?} In the fullest sense of that word, Man must encompass all of Creation: its lowest depths as surely as its highest heights. But the moral movement is upward, Without the upward moral and spiritual movement, all is chaos, an analogy for what will be more defined in Hades, which takes its meaning precisely from that upward moral movement. Without an upward movement, there is no ``Down'' either.

Modern man has at his beck and call all the primordial tradition of the past, scattered about like the body of Osiris, in the form of vast potential knowledge and development in many sciences which were not available to the ancient world in total form. What use are we making of those resources? And yet, Cologero has noted that ancient man (enviable in many ways compared to the modern version) was nonetheless degenerate\footnote{\url{https://www.gornahoor.net/?p=2577}}, that they were circumscribed by Fate just as we are. Everywhere and at all times, ``History is a crime'' (Nicholas Berdyaev). Those times of ignorance have past. The guardians and tutors have withdrawn, and the winter of the world has come, for there is a spiritual evolutionary test being administered, to see who can stand (in popular parlance, we might call this the ``Matrix''). Since ``now is the day of salvation'', this time period was the best time for you to appear.

\begin{quotex}
And the times of this ignorance God winked at; but now commandeth all men every where to repent: Because he hath appointed a day, in the which he will judge the world in righteousness by that man whom he hath ordained; whereof he hath given assurance unto all men, in that he hath raised him from the dead.
\flright{\textsc{Acts 17:30-31}}
\end{quotex}


There are ``cracks in the great wall'', as one author puts it.\footnote{\url{https://www.amazon.com/Cracks-Great-Wall-Traditional-Metaphysics/dp/1597310247}} For those who can read the signs of the Times, the Reign of Quantity\footnote{\url{https://en.wikipedia.org/wiki/The\_Reign\_of\_Quantity\_and\_the\_Signs\_of\_the\_Times}} symbolizes both the ``best of times'' and the ``worst of times''. Man will have to make progress, individually most of all, \emph{but he will be forced to it by the heavenly and the infernal powers, all the same, in an outer sense}. That is why Darwin, Marx, and Freud were given domination over the Modern Era. The motto of Freud\footnote{\url{https://en.wikipedia.org/wiki/Acheron}} was \emph{flectere si nequeo superos, Acheronta movebo}: `If I cannot bend the will of Heaven, I shall move Hell.''

This is why we inhabit, simultaneously, the most important of times (the present is always the scene of the struggle), and the least enviable of times (karmic debts have mounted). The Scriptures promise that God would not destroy the world with water again; the day of judgement will be a day of Fire . This Fire will either be experienced as Love, or as terrible Fire, depending upon the spiritual condition.

\begin{quotex}
The Lord is not slow about his promise as some count slowness, but is forbearing toward you, a not wishing that any should perish, but that all should reach repentance. But the day of the Lord will come like a thief, and then the heavens will pass away with a loud noise, and the elements will be dissolved with fire, and the earth and the works that are upon it will be burned up. Since all these things are thus to be dissolved, what sort of persons ought you to be in lives of holiness and godliness.
\flright{\textsc{2 Peter 3:9-11}}
\end{quotex}


We have behind us birth and baptism (Water) and in front of us Fire (the energetic state of Love that is destiny, and also the consequences for falling short). The entire metaphysics of Jacob Boehme\footnote{\url{http://www.jacobboehmeonline.com/}} explain the workings of the four elements in man and in Creation, which are combining to create and purify the spirit of Man, and they can be related to these ideas and patterns of God allowing (in the fallen state), Nature to urge and constrain itself to find a way out ``further in''. This perfect and theotic state\footnote{N Berdyaev (\url{http://www.krotov.info/library/02\_b/berdyaev/1941\_\_038\_0\_eng.htm}) terms man a ``theocosm'' and not just a ``microcosm''.} will not come without struggle and effort, and we will continue to see this conflict waged in the projected, externalized world, which groans for redemption. This is why Kukai\footnote{\url{https://plato.stanford.edu/entries/kukai/}} could say, ``The Mind and the World Co-Arise''.

Hercules will have more labors to accomplish. Defeating one enemy, and completing one labor, is not enough. He must work through the entire Zodiac circle\footnote{\url{http://www.bailey.it/files/Labours-of-Hercules.pdf}}, for this is the Cosmic Law. But in the process, He will be refined, until the day when ``heaven and earth kiss'', when the double oath of God, on earth as it is in heaven, is complete. Every time, every age, starts afresh in a sense -- God has set eternity in our hearts. The answer will not be found in one outer battle or triumph, which may just turn out to be a diversion or distraction. Do not foolishly ask ``why is this time so much worse than other times?'' The good old ways, the paths of the Lord, are found within you, for it is in the lines of the human heart that the meaning, the purpose, and the answer will be found. Along that path will you find the answers to all the other questions. Seek first the Kingdom of God, and His Righteousness, and all these things will be added unto you. It is a matter of winning the Greater War and the Inner Struggle. It always was. And it always will be, until it is done. Beware of seeking your true self -- you might find God. Beware of seeking God -- it may be that you will find your true self. This was vividly illustrated in the story of a Catholic martyr\footnote{\url{http://www.natcath.org/NCR\_Online/archives2/2007d/110907/110907a.htm}} during WWII, when the self and God no longer had need even for the Scriptures or outer forms.

\begin{quotex}
She said that Jochmann entered Jägerstätter's cell the night before the execution. He was the priest who said the prisoner refused to sign the document that could have spared his life. He said Jägerstätter had already received the last sacraments in the afternoon. Jochmann offered to bring Jägerstätter devotional reading material but the prisoner declined and he also declined to hear readings from scripture. According to the priest, Jägerstätter explained, ``I am completely bound in inner union with the Lord, and any reading would only interrupt my communication with my God.''
\end{quotex}


The labors of Hercules contain the keys to unlocking the fairy tale, and awakening the sleeping beauty, transforming Cinderella, or winning the favor of the Valkyries. This eighth labor is particularly notable for the devouring presence of the unintegrated parts of consciousness (the ``mares'' of the Anima\footnote{\url{https://www.blogger.com/null}}), because ``hell hath no fury like a woman scorned''. The overt tyrant or stereotypical, over-compensated ``male'' will not be able to control these forces, and in fact, not even a plain ``good man'' can control them either, hence the death of Hercules' companion. Only the presence of the true man on the quest can hope to quell and tame these forces, which later go on to serve usefully as bloodlines for great steeds for Alexander the Great.

To quote Cologero\footnote{\url{http://www.gornahoor.net/?p=7351}}:

\begin{quotex}
Man is self-sufficient, since in him the anima is the passive element. To give birth to the True Self, the anima must be pacified so that it can receive the imprint of a higher source. That is why Jung says that a physical woman is not absolutely necessary for the individuation process. Rather, she is a privation, so the energy of the anima exists in the man and can be claimed in an act of self-possession.
\end{quotex}


You may also notice that Hercules is not a sage or an intellectual: he is a warrior and an ascetic. We live in an era dominated by the Lunar influence of Intellect, and under its thrall. The age to come will not be one of Philosophy, Religion, or Science, but of Art, the Fourth Age, which is the unity of the preceding three. This is symbolized in John's Gospel, the fourth Gospel, which speaks the entire and full truth of the divinity of the Logos. But of the three kinds of men (natural, religious, and intellectual), it is the religious ``knights'' who will be most needed and called upon to undertake the artful transformation and balancing of man into the higher form. Natural or pagan man has no ``path'', and intellectuals today think only with their tongues. The men of ``heart'' cannot consent to either path, and are forced into transformation. You can find more teaching on this in Mouravieff's Gnosis\footnote{\url{https://www.bibliotecapleyades.net/esp\_autor\_mouravieff.htm}}, which is part of the basis of the Gornahoor web seminars started by Cologero.

Knowing things with the mind will help clear cobwebs away, but it isn't the same thing as beginning the journey. And it is interesting that Alice Bailey places the ``Eighth'' labor of Hercules, by rotation of the Zodiac, in the place of the First Labor. We could relate this to the First Arcanum of the Tarot, the Magician, which is the foundation for practicing the others. This would make sense, given that experientially knowing the Anima places a ``knight'' in a position of potency and power for accomplishing all of the subsequent Labors.

Follow your clues and commence the Quest.


\flrightit{Posted on 2018-05-29 by Logres}

